% Einheit 4 von 4 – Ausklammern (bank/terme/e4.jsonl, zusammenbau v0.9)
\einheitenkopf[e4]{Einheit 4 von 4 · Ausklammern}
\zweigzeile{Hier lernst du, auszuklammern – gemeinsamer Zahlfaktor, gemeinsame Variable, Probe durch Ausmultiplizieren · ab Kl. 8 (am Gymnasium ab Kl. 7) · keine P10-Aufgabe · baut auf: Addieren und Subtrahieren negativer Zahlen, Punkt vor Strich mit Zahlen}
\setcounter{aufgabe}{34}


% erkennung: terme-e4-k1-s0-v1, terme-e4-k1-s0-v2, terme-e4-k1-s0-v3, terme-e4-k1-s0-v4
\begin{kbaufgabe}{Ich erkenne, welcher Faktor in jedem Glied steckt.}
\begin{kbblock}
\kbanweisung{Was steckt in jedem Glied von …? Kreise den Faktor ein, der in beiden Gliedern steckt. Klammere nichts aus.}
\lbpaar{\kbteil{a)}{$8x + 12$}{}
}{\kbteil{b)}{$9y^2 + 2y$}{}}
\end{kbblock}
\begin{kbblock}
\kbteil{c)}{Was steckt in jedem Glied von $6a^2 + 10a$? Kreise ein, was in beiden Gliedern steckt: Zahl, Variable oder beides. Klammere nichts aus.}{}
\end{kbblock}
\begin{kbblock}
\kbteil{d)}{Was steckt in jedem Glied von $9x - 6 + 15x^2$? Kreise ein, was in allen drei Gliedern steckt. Klammere nichts aus.}{}
\end{kbblock}
\end{kbaufgabe}


% vorstufe: terme-e4-k2-s-1-v1, terme-e4-k2-s-1-v2, terme-e4-k2-s-1-v3, terme-e4-k2-s-1-v4
\begin{kbaufgabe}{Ich kann jedes Glied als Produkt mit einem Faktor schreiben.}
\begin{kbblock}
\lbpaar{\kbteil{a)}{Schreibe jedes Glied von $8x + 20$ als Produkt mit dem Faktor 4.}{}
\kbantwort{$4 \cdot$ \kbfeld $+ 4 \cdot$ \kbfeld}
}{\kbteil{b)}{Schreibe jedes Glied von $6a + 14$ als Produkt mit dem Faktor 2.}{}
\kbantwort{$2 \cdot$ \kbfeld $+ 2 \cdot$ \kbfeld}}
\end{kbblock}
\begin{kbblock}
\lbpaar{\kbteil{c)}{Schreibe jedes Glied von $10y + 15$ als Produkt mit dem Faktor 5.}{}
\kbantwort{$5 \cdot$ \kbfeld $+ 5 \cdot$ \kbfeld}
}{\kbteil{d)}{Schreibe jedes Glied von $9b + 6$ als Produkt mit dem Faktor 3.}{}
\kbantwort{$3 \cdot$ \kbfeld $+ 3 \cdot$ \kbfeld}}
\end{kbblock}
\end{kbaufgabe}


% vorstufe: terme-e4-k2-s0-v1, terme-e4-k2-s0-v2, terme-e4-k2-s0-v3, terme-e4-k2-s0-v4
\begin{kbaufgabe}{Ich kann einen vorgegebenen Faktor ausklammern.}
\begin{kbblock}
\lbpaar{\kbteil{a)}{Klammere den Faktor 3 aus: $9x + 21$.}{}
\kbantwort{$3 \cdot ($ \kbfeld $+$ \kbfeld $)$}
}{\kbteil{b)}{Klammere den Faktor 2 aus: $4y + 10$.}{}
\kbantwort{$2 \cdot ($ \kbfeld $+$ \kbfeld $)$}}
\end{kbblock}
\begin{kbblock}
\lbpaar{\kbteil{c)}{Klammere den Faktor 5 aus: $10a + 35$.}{}
\kbantwort{$5 \cdot ($ \kbfeld $+$ \kbfeld $)$}
}{\kbteil{d)}{Klammere den Faktor 7 aus: $14b + 21$.}{}
\kbantwort{$7 \cdot ($ \kbfeld $+$ \kbfeld $)$}}
\end{kbblock}
\end{kbaufgabe}


% leiter: terme-e4-k2-s1-v1, terme-e4-k2-s1-v2, terme-e4-k2-s1-v3, terme-e4-k2-s1-v4, terme-e4-k2-s2-v1, terme-e4-k2-s3-v1, terme-e4-k2-s4-v1, terme-e4-k2-s5-v1, terme-e4-k2-s6-v1, terme-e4-k2-s7-v1, terme-e4-k2-s8-v1, terme-e4-k2-s9-v1
\begin{kbaufgabe}{Ich kann ausklammern.}
\begin{kbblock}
\kbanweisung{Klammere den größten gemeinsamen Faktor aus.}
\lbpaar{\kbteil{a)}{$5x + 10 =$ \lbfeld}{}
}{\kbteil{b)}{$5x + 25 =$ \lbfeld}{}}
\end{kbblock}
\begin{kbblock}
\lbpaar{\kbteil{c)}{$5x + 30 =$ \lbfeld}{}
}{\kbteil{d)}{$5x + 35 =$ \lbfeld}{}}
\end{kbblock}
\begin{kbblock}
\lbpaar{\kbteil{e)}{$x^2 + 5x =$ \lbfeld}{}
}{\kbteil{f)}{$6x^2 + 15x =$ \lbfeld}{}}
\end{kbblock}
\begin{kbblock}
\lbpaar{\kbteil{g)}{$6x + 6 =$ \lbfeld}{}
}{\kbteil{h)}{$3x^2 + 6x + 15 =$ \lbfeld}{}}
\end{kbblock}
\begin{kbblock}
\kbanweisung{Klammere aus und mache die Probe.}
\kbteil{i)}{$6x + 30$}{}
\item[]\kbraum{2}
\end{kbblock}
\begin{kbblock}
\kbteil{j)}{Zwei Rechtecke liegen nebeneinander. Beide sind 4 cm hoch, das eine ist $x$ cm breit, das andere 6 cm. Sie haben die Höhe als gemeinsame Seite. Gib den Flächeninhalt der ganzen Figur in cm² als Summe der Teilflächen und als Produkt mit Klammer an.}{}
\item[]\kbraum{2}
\end{kbblock}
\begin{kbblock}
\kbteil{k)}{Eine Hose kostet 40 €, ein Hemd 25 €. Es gibt 10 \% Rabatt. Anna rechnet den Rabatt mit $0{,}1 \cdot 40 + 0{,}1 \cdot 25$, Ben mit $0{,}1 \cdot (40 + 25)$. Begründe, ob beide Terme denselben Rabatt ergeben, und berechne den Rabatt.}{}
\item[]\kbraum{2}
\end{kbblock}
\begin{kbblock}
\kbteil{l)}{Nora soll im Term $5x + 20$ den größten gemeinsamen Faktor ausklammern. Sie rechnet so: \rechnung{5x + 20 &= 5 \cdot (x + 20)} Finde den Fehler und rechne richtig.}{}
\item[]\kbraum{2}
\end{kbblock}
\end{kbaufgabe}


% pruefung: terme-e4-k2-s10-v1
\begin{kbaufgabe}{Ich kann auch Terme mit zwei Variablen so weit wie möglich ausklammern.}
\begin{kbblock}
\kbanweisung{Klammere so weit wie möglich aus und mache die Probe.}
\kbteil{}{$3a^2 + 6ab + 9a$}{}
\item[]\kbraum{2}
\end{kbblock}
\end{kbaufgabe}


% pflicht: terme-e4-k3-s1-v1, terme-e4-k3-s1-v2, terme-e4-k3-s1-v3
\begin{kbaufgabe}{Ich finde den Fehler beim Ausklammern.}
\begin{kbblock}
\kbteil{a)}{Finn soll im Term $20x - 8$ den größten gemeinsamen Faktor ausklammern. Er rechnet so: \rechnung{20x - 8 &= 4 \cdot (5x + 2)} Finde den Fehler und rechne richtig.}{}
\item[]\kbraum{2}
\end{kbblock}
\begin{kbblock}
\kbteil{b)}{Sara soll im Term $10a - 15$ den größten gemeinsamen Faktor ausklammern. Sie rechnet so: \rechnung{10a - 15 &= 5 \cdot (2a - 3)} Prüfe, ob Sara richtig gerechnet hat.}{}
\item[]\kbraum{2}
\end{kbblock}
\begin{kbblock}
\kbteil{c)}{Max hat in vier Termen ausgeklammert. Welche Ergebnisse können nicht stimmen? Begründe, ohne genau zu rechnen.}{}
\kbfrage{(1) $6a + 18 = 6 \cdot (a + 3)$ \\ (2) $10x - 25 = 5 \cdot (2x + 5)$ \\ (3) $7y + 7 = 7 \cdot y$ \\ (4) $9b - 36 = 9 \cdot (b - 4)$}
\item[]\kbraum{2}
\end{kbblock}
\end{kbaufgabe}


% pflicht: terme-e4-k3-s2-v1, terme-e4-k3-s2-v2, terme-e4-k3-s2-v3
\begin{kbaufgabe}{Ich kann begründen, ob richtig ausgeklammert wurde.}
\begin{kbblock}
\kbteil{a)}{Begründe, warum die Terme $6a + 42$ und $6 \cdot (a + 7)$ für jede Zahl $a$ denselben Wert haben. Prüfe es auch durch Einsetzen von $a = 2$.}{}
\item[]\kbraum{3}
\end{kbblock}
\begin{kbblock}
\kbteil{b)}{Entscheide bei jeder Aussage, ob sie wahr oder falsch ist. Begründe, ohne genau zu rechnen.}{}
\kbfrage{(1) Aus jedem Term mit zwei Gliedern kann man einen Faktor größer als 1 ausklammern. \\ (2) Multipliziert man nach dem Ausklammern wieder aus, erhält man immer den Anfangsterm. \\ (3) Klammert man aus zwei Gliedern aus, stehen in der Klammer nie weniger als zwei Glieder.}
\item[]\kbraum{3}
\end{kbblock}
\begin{kbblock}
\kbteil{c)}{Lina sagt: „Bei $8x - 8$ klammere ich 8 aus, in der Klammer bleibt dann nur $x$.“ Begründe, ob Lina recht hat.}{}
\item[]\kbraum{3}
\end{kbblock}
\end{kbaufgabe}


% pflicht: terme-e4-k3-s3-v1, terme-e4-k3-s3-v2, terme-e4-k3-s3-v3
\begin{kbaufgabe}{Ich kann einen Term mit Klammer in Worten oder als Figur beschreiben.}
\begin{kbblock}
\kbteil{a)}{Der Term $6 \cdot (b + 3)$ gibt den Flächeninhalt einer Figur aus zwei Rechtecken in cm² an. Beschreibe die Figur in Worten, ohne sie zu zeichnen.}{}
\item[]\kbraum{2}
\end{kbblock}
\begin{kbblock}
\kbteil{b)}{Man nimmt die Summe aus einer Zahl $x$ und 9 mal 4. Schreibe als Term mit Klammer. Schreibe den Term dann ohne Klammer.}{}
\item[]\kbraum{2}
\end{kbblock}
\begin{kbblock}
\kbteil{c)}{Beschreibe den Term $3 \cdot (y + 5)$ in Worten, ohne ihn umzuformen.}{}
\item[]\kbraum{2}
\end{kbblock}
\end{kbaufgabe}


% pflicht: terme-e4-k3-s4-v1, terme-e4-k3-s4-v2, terme-e4-k3-s4-v3
\begin{kbaufgabe}{Ich kann mit Ausklammern Aufgaben aus dem Alltag lösen.}
\begin{kbblock}
\kbteil{a)}{Ein Sportverein kauft 15 Trikots zu je 28 € und 15 Hosen zu je 17 €. Schreibe die Kosten als Term mit Klammer. Der Verein hat 700 €. Reicht das Geld?}{}
\item[]\kbraum{2}
\end{kbblock}
\begin{kbblock}
\kbteil{b)}{Frau Kaya kauft für jedes ihrer 3 Kinder eine Jacke zu 45 € und ein Paar Schuhe zu 38 €. Stelle einen Term mit Klammer für den Gesamtpreis auf. Wie viel Euro kostet der Einkauf?}{}
\item[]\kbraum{2}
\end{kbblock}
\begin{kbblock}
\kbteil{c)}{Zwei rechteckige Beete liegen nebeneinander. Beide sind 4 m breit, das eine ist 7,5 m lang, das andere 5,5 m. Ein Sack Rasendünger reicht für 50 m². Reicht ein Sack für beide Beete?}{}
\item[]\kbraum{2}
\end{kbblock}
\end{kbaufgabe}
